

 In order to formulate our main results on the scattering problem we
need to introduce ap\-propriate Hilbert spaces and auxiliary norms.
By $\vert n\rangle$, $n=(n_1,\dots,n_k)\in\mathbb N_0^d$ we shall denote the
standard orthonormal basis for $L^2(\mathbb R^d)$ consisting of
eigenstates for the $d$-dimensional harmonic oscillator, and by
$(\phi_{mn})$ we denote the matrix representing a bounded operator
$\phi$ with respect to this basis, that is
\[
\phi = \sum_{m,n}\phi_{mn}\vert n\rangle\langle m|,\qquad \phi_{mn}:=\langle n|\phi|m\rangle .
\]
Let
\[
b_{mn} := 1+|m-n| ,\qquad m,n\in \mathbb N_0^d ,
\]
where $|\cdot|$ denotes the Euclidean norm in $\mathbb R^d$. For
exponents $p\geq 1$, $\alpha\in\mathbb R$ and $\phi$ as above, we def\/ine
the norms $\Vert\cdot \Vert_{p,\alpha}$  by
\begin{gather}
\label{alphanorms}
\Vert\phi\Vert_{p,\alpha}^p := \sum_{m,n}b_{mn}^{\alpha p}|\phi_{mn}|^p ,
\quad p<\infty, \qquad\mbox{and}\qquad
\Vert\phi\Vert_{\infty,\alpha}:=\sup_{m,n}\{b_{mn}^\alpha|\phi_{mn}|\} ,
\end{gather}
and we let ${\cal L}_{p,\alpha}$ denote the space of operators
on $L^2(\mathbb R^d)$ for which
$\Vert\cdot \Vert_{p,\alpha}$ is f\/inite. We note that
$\mathcal H_{\alpha}:=\mathcal L_{2,\alpha}$ is a  Hilbert space with
scalar product
\[
\langle\phi,\psi\rangle_\alpha:=\sum_{m,n}b_{mn}^{2\alpha} \overline{\phi_{mn}}
 \psi_{mn} ,
\]
and, in particular, ${\cal H}_{0}$ equals  the space $\cal H$ of
Hilbert--Schmidt operators on $L^2(\mathbb R^d)$ with $\Vert\cdot\Vert_{2,0} =
\Vert\cdot\Vert_2$.
Clearly, the operator $U_\alpha:\mathcal H_\alpha\to \mathcal H$ def\/ined by
\begin{gather}\label{Ualpha}
(U_\alpha\phi)_{mn}=
b^{\alpha}_{mn}\phi_{mn} ,\qquad \phi\in\mathcal H_\alpha ,
\end{gather}
is unitary, and the operator
${\bf\Delta}_\alpha := U_\alpha^*\, {\bf \Delta}\,U_\alpha$, with domain
$D({\bf\Delta}_\alpha):=U_\alpha^* D({\bf \Delta})$,
 is self-adjoint on $\mathcal H_\alpha$.

We use the notation $\Vert\cdot\Vert_{\rm op}$ for
the standard operator norm and def\/ine the norm $\Vert\cdot\Vert_a$ by
\[
\Vert\phi\Vert_a := \Vert\widetilde\phi\Vert_{\rm op} ,
\]
whenever the operator
\[
 \widetilde\phi := \sum_{m,n}|\phi_{mn}|\, |n\rangle\langle m|
\]
 is bounded.


With this notations we have the following decay estimate for the free
propagation:
\begin{theorem}\label{th2} For $\alpha\geq 0$ and  $\beta > d$ there exists a constant
$c_\beta>0$ such that
\begin{gather}\label{freedecay}
\Vert\, e^{-it{\bf\Delta}_\alpha}\phi \Vert_a  \leq  c_{\beta}(1+
|t|)^{-\frac{d}{2}} \Vert\phi\Vert_{1,\beta} ,
\qquad \phi\in {\cal H}_\alpha\cap {\cal L}_{1,\beta} .
\end{gather}
\end{theorem}

This estimate should be compared to the standard one used for the
classical nonlinear wave equations with the $L^\infty$-norm on the
left and a $L^1$-Sobolev norm on the right, applied to
functions of $2d$ space variables. In this case, one obtains rather
trivially a decay exponent~$d$ instead of~$d/2$. However, those norms
are not well behaved with respect to the $\star$-product and cannot be
used for our purposes. Instead, we have to work a bit harder to establish
(\ref{freedecay}) using uniform estimates on the classical Jacobi
polynomials. This is accomplished in Section~\ref{secth2}.



Our last main result concerns the existence of wave operators
def\/ined on a scattering subspace $\Sigma_\alpha\subset \mathcal H_\alpha$. In order to def\/ine
the latter we introduce, for f\/ixed $\alpha\geq 0$, the scattering norm
of $\phi\in{\cal H}_\alpha$ by
\begin{gather*}%\label{scatnorm}
|||\phi|||_{\alpha}:=
\Vert\phi\Vert_{2,\alpha}+\sup_{t\in\mathbb R} |t|^{\frac d2}
\Vert e^{-it{\bf\Delta}_\alpha}\phi\Vert_a
\end{gather*}
and set
\begin{gather*}
%\label{sigmascat}
\Sigma_{\alpha}:=\left\{\phi\in\mathcal H_\alpha\,|\;|||\phi|||_{\alpha}<\infty\right\}\,.
\end{gather*}
From \eqref{alphanorms}, we immediately see that
$\mathcal L_{p,\alpha}\subset\mathcal L_{q,\beta}$, if $p\leq q$ and
$\alpha p\geq\beta q$, so that Theorem \ref{th2} gives, in particular, the
inclusion $\mathcal L_{1,\beta}\subset\Sigma_\alpha$ if $\beta> \max\{d, 2\alpha\}$.

The relevant integral equations to solve in a scattering situation
correspond to initial/f\/inal data $\phi_\pm$ at $t=\pm\infty$ and
take the form
\begin{gather}
\label{NLSch4-}
\phi(t)=e^{-it{{\bf\Delta}_\alpha}}\phi_--i\int_{-\infty}^t
e^{i{(s-t){\bf\Delta}_\alpha}}\phi(s) F\big(|\phi(s)|^2\big)\,ds ,
\end{gather}
and
\begin{gather}
\label{NLSch4+}
\phi(t)=e^{-it{{\bf\Delta}_\alpha}}\phi_++i\int^{+\infty}_t
e^{i{(s-t){\bf\Delta}_\alpha}}\phi(s) F\big(|\phi(s)|^2\big)\,ds ,
\end{gather}
respectively, where $F$ has been redef\/ined to include $\theta$.
 {\it We assume here for simplicity that the polynomial $F$
has no constant term}, since such a term could trivially be incorporated by
subtracting it from ${\bf\Delta}_\alpha$ in the preceding formulas.   The
f\/irst problem
then is to determine spaces of initial/f\/inal data $\phi_\pm$, such that the
equations above have a unique global solution and such that these
solutions behave as free solutions for $t\to \pm\infty$.
To this end we establish the following in Section~\ref{secth3}.

\begin{theorem}\label{th3}
Let $\alpha>2d$ and assume that the polynomial $F$ has no constant
term and, in addition, no linear term
if $d=1$ or $d=2$.

Then there exists $\delta>0$ such that for every
 $\phi_\pm\in\Sigma_{\alpha}$ with $|||\phi_\pm|||_{\alpha}<\delta$, the
equations~\eqref{NLSch4-} and~\eqref{NLSch4+} have unique globally
defined continuous solutions $\phi^\pm:\mathbb R\to \Sigma_{\alpha}$ fulfilling
\begin{gather}\label{solsmall1}
\Vert\phi^\pm(t)- e^{-it{{\bf\Delta}_\alpha}}\phi_\pm\Vert_{2,\alpha}  \to
0,\qquad\mbox{for}\ \ t\to\pm\infty .
\end{gather}
If, furthermore, $F$ has no linear term, then we have
\begin{gather}\label{solsmall2}
 |||e^{it{{\bf\Delta}_\alpha}}\phi^\pm(t)- \phi_\pm |||_{\alpha}  \to
0,\qquad\mbox{for}\  \ t\to\pm\infty .
\end{gather}
\end{theorem}
